\clearpage
\section[Conclusion and Future Scope]{Chapter -- 04: Conclusion and Future Scope}

This chapter summarises what the results in Chapter~3 establish about profiled,
machine-learning side-channel attacks on masked ASCAD traces, and identifies concrete
extensions of this work.

\subsection{Conclusion}

This project rebuilt the profiled side-channel attack of Poudel and
Rahimi~\cite{poudel2025} end to end -- data loading, preprocessing, four model
families, and a Key Rank / Guessing Entropy evaluation -- and checked which of its
qualitative claims survive that reconstruction. The classical-versus-deep split
reproduces cleanly: neither the random forest nor the SVM ever rank the true key
byte, while the compact CNN and the 1-D ResNet both recover it on the fixed-key set,
and the ResNet does so more efficiently than the CNN on the harder variable-key set,
which is the paper's central qualitative point. Two extensions beyond the paper
produced genuinely new findings. First, a multi-task ResNet that also predicts the
mask shares cuts the fixed-key trace requirement roughly threefold (870 to 270) by
breaking the early training plateau sooner, but is actively worse on the
variable-key set, where that plateau is not the limiting factor -- a useful reminder
that an auxiliary training objective only helps while optimisation, not capacity, is
the bottleneck. Second, the ResNet tolerates trace desynchronisation that defeats the
compact CNN outright, matching the wider literature's finding that deeper residual
networks are more shift-tolerant than compact convolutional networks. A
shuffled-label control confirmed the pipeline depends on genuine trace-label
association rather than a label-distribution artefact. The efficiency figures
reported in the wider literature ($\sim$35 traces) were not matched -- the best
result here is roughly 8$\times$ that -- and closing this gap is the main open target
for further work.

\subsection{Future Scope}

\begin{itemize}
\item Running a proper multi-seed hyperparameter search on the variable-key ResNet
      to separate genuine efficiency gains from seed luck, and to close the gap to
      the literature's $\sim$35-trace figure.
\item Adding shift augmentation during training to close the remaining gap between
      aligned and desynchronised ResNet performance.
\item Acquiring the raw ATMega8515 trace set and implementing a per-byte
      re-windowing step, to attempt full 16-byte key recovery rather than byte 2
      alone.
\item Extending the multi-task objective's early-plateau benefit to the
      variable-key set, for example by weighting the auxiliary heads down once the
      main task starts improving.
\item Running the desynchronisation and multi-task experiments across multiple
      seeds to replace single-run figures with mean $\pm$ std results.
\item Applying the same Key Rank benchmark methodology to other masked SCA datasets
      or leakage models (e.g.\ the Hamming-weight model, or other AES
      implementations) as they become available.
\end{itemize}

\textbf{Result:} The project's objectives (Chapter~1) were fully met: a working,
end-to-end key-recovery pipeline was implemented (Chapter~2), evaluated across
classical and deep model families with three extensions beyond the original paper
(Chapter~3), and shown to reproduce the paper's central qualitative claims while
surfacing a genuine limitation -- the gap to the literature's best reported
trace-efficiency figures -- that further tuning would need to close.
